% =====================================================================
%  Приложение A. Сводная таблица обозначений
%
%  ЖАНР (STYLE.md): учебная литература. Свод символов, введённых в §0-§10,
%  сгруппированный по темам в порядке появления.
% =====================================================================
\section{Сводная таблица обозначений}
\label{sec:notation}
\markboth{Приложение. Сводная таблица обозначений}{}

Ниже собраны все обозначения курса в порядке их первого появления, сгруппированные по
тематическим блокам. Курс не вводит собственной экзотической нотации~--- везде, где это возможно,
используются общепринятые в оптике и статистике символы; отступления оговорены отдельно.

% ---------------------------------------------------------------------
\subsection*{Электродинамика и поле (§0)}
\begin{center}
\footnotesize
\setlength{\tabcolsep}{6pt}
\begin{tabular}{@{}p{0.16\linewidth}p{0.84\linewidth}@{}}
\toprule
\textbf{Символ} & \textbf{Значение} \\
\midrule
$\nu$, $\omega$, $\lambda$, $k$ & частота (Гц или ТГц), круговая частота $\omega=2\pi\nu$, длина
волны $\lambda=c/\nu$, волновое число $k=2\pi/\lambda$ \\
$A$, $\varphi_0$ & амплитуда и начальная фаза вещественной волны $E(z,t)=A\cos(\omega t-kz+\varphi_0)$ \\
$\Efield=A e^{\jj\varphi_0}$ & комплексная амплитуда (макрос \texttt{\textbackslash Efield}) \\
$\jj$ & мнимая единица в инженерной записи (макрос \texttt{\textbackslash jj}), $\jj^2=-1$ \\
$\tau$ (общее) & временная задержка, извлекаемая из наклона фазы по частоте, $\delta(\nu)=\delta_0+2\pi\nu\tau$ \\
$U(\theta)$ & наблюдаемая энергия (интенсивность) на угле~$\theta$ \\
\bottomrule
\end{tabular}
\end{center}

\subsection*{Геометрия решётки (§0, §2)}
\begin{center}
\footnotesize
\setlength{\tabcolsep}{6pt}
\begin{tabular}{@{}p{0.16\linewidth}p{0.84\linewidth}@{}}
\toprule
$D$, $P$ & диаметр проволоки и период решётки \\
$D/P$ & параметр плотности решётки (доля периода, занятая металлом) \\
$\Deff$ & эффективный диаметр (макрос \texttt{\textbackslash Deff}) — значение $D$, которое
даёт согласие модели с измерением \\
$\Dphys$ & физический диаметр (макрос \texttt{\textbackslash Dphys}), измеренный микроскопией \\
$\tperp(\nu)$, $\tpar(\nu)$ & комплексные коэффициенты пропускания канала пропускания и канала
подавления (макросы \texttt{\textbackslash tperp}, \texttt{\textbackslash tpar}) \\
$\theta$ & угол ориентации оси пропускания элемента в лабораторных осях; провода лежат под
$\theta+90^\circ$ \\
$\theta_0$ & офсет угла (ноль шкалы поворотного столика относительно физической оси элемента) \\
\bottomrule
\end{tabular}
\end{center}

\subsection*{Матрицы Джонса (§1)}
\begin{center}
\footnotesize
\setlength{\tabcolsep}{6pt}
\begin{tabular}{@{}p{0.16\linewidth}p{0.84\linewidth}@{}}
\toprule
$\mathbf{E}=(\Efield_x,\Efield_y)^{\mathsf T}$ & вектор Джонса (состояние поляризации) \\
$\mathbf{J}$ & матрица Джонса оптического элемента, $2\times2$, комплексная \\
$\Rot{\theta}$ & матрица поворота на угол~$\theta$ (макрос \texttt{\textbackslash Rot}) \\
$\eta=|\tpar/\tperp|$ & просачивание (амплитудное отношение каналов) \\
$\delta=\arg(\tpar/\tperp)$ & фаза просачивания; $\delta(\nu)=\delta_0+2\pi\nu\tleak$ \\
$\tleak$ & групповая задержка просачивания (макрос \texttt{\textbackslash tleak}) \\
$\mathrm{ER}=1/\eta^2$ & коэффициент экстинкции (по мощности) \\
$\mathbf{d}$ & вектор чувствительности приёмника \\
\bottomrule
\end{tabular}
\end{center}

\subsection*{Модель Бланко (§2)}
\begin{center}
\footnotesize
\setlength{\tabcolsep}{6pt}
\begin{tabular}{@{}p{0.16\linewidth}p{0.84\linewidth}@{}}
\toprule
$A_1$, $A_2$ & безразмерные функции геометрии, входящие в оба канала (не путать с амплитудами
гармоник $A_2,A_4$ §4~--- контекстно различимые обозначения, одноимённые по историческим причинам
источника) \\
$C_m=\sqrt{m^2-(P/\lambda)^2}$ & слагаемое суммы по дифракционным порядкам \\
$Z_a,\ldots,Z_d$ & нормированные реактивные импедансы эквивалентной схемы \\
\bottomrule
\end{tabular}
\end{center}

\subsection*{Материальные поправки (§3)}
\begin{center}
\footnotesize
\setlength{\tabcolsep}{6pt}
\begin{tabular}{@{}p{0.16\linewidth}p{0.84\linewidth}@{}}
\toprule
$\sigma(\omega)$, $\tau_e$ & проводимость металла по модели Друде и время релаксации электрона \\
$Z_s$, $Z_0$ & поверхностный импеданс металла и импеданс вакуума \\
$\gamma$ & показатель степенного члена потерь \\
$L$ (в §3) & амплитудный коэффициент степенных потерь, дБ на единицу частотной шкалы (контекстное
обозначение, не путать с правдоподобием $L$ §8) \\
\bottomrule
\end{tabular}
\end{center}

\subsection*{Линейная аппроксимация и корреляция (§4-§5)}
\begin{center}
\footnotesize
\setlength{\tabcolsep}{6pt}
\begin{tabular}{@{}p{0.16\linewidth}p{0.84\linewidth}@{}}
\toprule
$a_0,a_1,b_1$ & коэффициенты учебной 3-параметрической формы \\
$c_0,a_2,b_2,a_4,b_4$ & коэффициенты точной 5-параметрической формы; $c_2=\sqrt{a_2^2+b_2^2}$,
$c_4$~--- аналогично для 4-й гармоники \\
$X$ & матрица плана (столбцы~--- значения базисных функций на измеренных углах) \\
$W=\operatorname{diag}(1/\sigma_i^2)$ & диагональная матрица весов \\
$\operatorname{Cov}(\mathbf{p})$ & ковариационная матрица оценки параметров \\
$\cond(\cdot)$ & число обусловленности матрицы (макрос-оператор \texttt{\textbackslash cond}) \\
\bottomrule
\end{tabular}
\end{center}

\subsection*{Нелинейное оценивание (§6-§7)}
\begin{center}
\footnotesize
\setlength{\tabcolsep}{6pt}
\begin{tabular}{@{}p{0.16\linewidth}p{0.84\linewidth}@{}}
\toprule
$\mathbf{p}$ & вектор параметров модели \\
$\chi^2(\mathbf{p})$ & взвешенная сумма квадратов невязок (целевая функция) \\
$\eta_0$ & амплитуда просачивания как свободный параметр подгонки (частотная зависимость~---
$\eta(\nu)$ или $\eta_{\exp}$) \\
$\varepsilon_{\rm cross}$ & комплексный перекрёстный член тракта, не зависящий от угла \\
$\phi_2$, $\psi$ & разъюстировка оси второго элемента тракта и наклон излучателя \\
$\argmin$ & оператор минимизации по параметру (макрос-оператор \texttt{\textbackslash argmin}) \\
\bottomrule
\end{tabular}
\end{center}

\subsection*{Правдоподобие и информационный критерий (§8)}
\begin{center}
\footnotesize
\setlength{\tabcolsep}{6pt}
\begin{tabular}{@{}p{0.16\linewidth}p{0.84\linewidth}@{}}
\toprule
$L(\mathbf{p})$ & функция правдоподобия (контекстное обозначение, не путать с амплитудным
коэффициентом потерь $L$ §3) \\
$k$ (в §8) & число свободных параметров модели (контекстное обозначение, не путать с волновым
числом $k$ §0) \\
$\AIC=-2\ln L+2k$ & информационный критерий Акаике (макрос-оператор \texttt{\textbackslash AIC}) \\
$\dAIC$ & разность информационных критериев двух моделей (макрос \texttt{\textbackslash dAIC}) \\
\bottomrule
\end{tabular}
\end{center}

\subsection*{Идентифицируемость (§9)}
\begin{center}
\footnotesize
\setlength{\tabcolsep}{6pt}
\begin{tabular}{@{}p{0.16\linewidth}p{0.84\linewidth}@{}}
\toprule
$\mathcal{I}(\mathbf{p}^\star)$ & матрица информации Фишера в точке минимума \\
$J$ (в §9) & якобиан модели, $J_{i\alpha}=\partial f/\partial p_\alpha$ (контекстное
обозначение, не путать с матрицей Джонса $\mathbf{J}$ §1, обозначаемой полужирным) \\
\bottomrule
\end{tabular}
\end{center}

\begin{warning}[: о повторном использовании символов]
Курс, как и большинство прикладных текстов, не может избежать полностью повторного использования
букв разными разделами математики: $L$ означает то амплитуду потерь (§3), то правдоподобие (§8);
$k$~--- то волновое число (§0), то число параметров (§8); $J$~--- то матрицу Джонса (§1, полужирным
$\mathbf{J}$), то якобиан (§9). В каждом случае значение однозначно определяется контекстом
параграфа; там, где два обозначения могли бы столкнуться в одной формуле, курс явно оговаривает,
какое из них имеется в виду.
\end{warning}

\clearpage
